% Part B paper (Parallel Programming): IEEE conference format, 8 pages.
% Build: see README.md. Numbers come from generated/numbers.tex (paper/numbers.py); never type a number
% that a result file holds. Red [pending: ...] notes are the outline, to be replaced by prose.
\documentclass[conference]{IEEEtran}
\usepackage{cite}
\usepackage{amsmath,amssymb,amsfonts}
\usepackage{graphicx}
\usepackage{textcomp}
\usepackage{booktabs}
\usepackage{url}
\usepackage{xcolor}
% No character protrusion: it pushes line-end punctuation into the column gutter, which submission
% checkers measure (as on the ICEE 2026 paper).
\usepackage[protrusion=false]{microtype}

\graphicspath{{figures/}}
% Generated by paper/numbers.py from committed result files. Do not edit.
\newcommand{\PeakTflops}{7.46}  % results/cuda/peak.json
\newcommand{\PeakGhz}{1.9}  % results/cuda/peak.json
\newcommand{\DramGBs}{188}  % results/cuda/peak.json
\newcommand{\LaunchUs}{1.8}  % results/cuda/peak.json
\newcommand{\LaunchSyncUs}{4.6}  % results/cuda/peak.json
\newcommand{\PredVoneLo}{1.98}  % results/cuda/predicted_times.json
\newcommand{\PredVoneHi}{4.86}  % results/cuda/predicted_times.json
\newcommand{\PredVtwoLo}{1.16}  % results/cuda/predicted_times.json
\newcommand{\PredVtwoHi}{1.62}  % results/cuda/predicted_times.json
\newcommand{\PredVthreeLo}{0.67}  % results/cuda/predicted_times.json
\newcommand{\PredVthreeHi}{1.24}  % results/cuda/predicted_times.json
\newcommand{\PredBoneLo}{0.51}  % results/cuda/predicted_times.json
\newcommand{\PredBoneHi}{1.78}  % results/cuda/predicted_times.json
\newcommand{\PinMhz}{1395}  % results/attention/clock1395/summary.json
\newcommand{\PinVoneMs}{26.06}  % results/attention/clock1395/summary.json
\newcommand{\PinVonePredLo}{2.50}  % results/attention/clock1395/summary.json
\newcommand{\PinVonePredHi}{6.42}  % results/attention/clock1395/summary.json
\newcommand{\PinVtwoMs}{1.92}  % results/attention/clock1395/summary.json
\newcommand{\PinVtwoPredLo}{1.38}  % results/attention/clock1395/summary.json
\newcommand{\PinVtwoPredHi}{2.01}  % results/attention/clock1395/summary.json
\newcommand{\PinVthreeMs}{1.77}  % results/attention/clock1395/summary.json
\newcommand{\PinVthreePredLo}{0.91}  % results/attention/clock1395/summary.json
\newcommand{\PinVthreePredHi}{1.68}  % results/attention/clock1395/summary.json
\newcommand{\PinBmathMs}{3.08}  % results/attention/clock1395/summary.json
\newcommand{\PinBmathPredLo}{0.69}  % results/attention/clock1395/summary.json
\newcommand{\PinBmathPredHi}{2.22}  % results/attention/clock1395/summary.json
\newcommand{\PinBeffMs}{2.22}  % results/attention/clock1395/summary.json
\newcommand{\PinBeffPredLo}{0.69}  % results/attention/clock1395/summary.json
\newcommand{\PinBeffPredHi}{2.22}  % results/attention/clock1395/summary.json
\newcommand{\PinSavingMs}{0.14}  % results/attention/clock1395/summary.json
\newcommand{\PinSavingPctOfNxN}{13}  % results/attention/clock1395/summary.json
\newcommand{\PinVthreeOverVtwoPct}{7.5}  % results/attention/clock1395/summary.json
\newcommand{\DepVoneMs}{18.96}  % results/attention/v1.json
\newcommand{\DepVoneGhz}{1.95}  % results/attention/v1.json
\newcommand{\DepVtwoMs}{1.82}  % results/attention/v2.json
\newcommand{\DepVtwoGhz}{1.51}  % results/attention/v2.json
\newcommand{\DepVthreeMs}{1.42}  % results/attention/v3.json
\newcommand{\DepVthreeGhz}{1.78}  % results/attention/v3.json
\newcommand{\DepBmathMs}{2.84}  % results/attention/b1.json
\newcommand{\DepBmathGhz}{1.69}  % results/attention/b1.json
\newcommand{\DepBeffMs}{1.90}  % results/attention/b1.json
\newcommand{\DepBeffGhz}{1.69}  % results/attention/b1.json
\newcommand{\VthreeBigNMs}{118}  % results/attention/v3.json
\newcommand{\FixedVoneUs}{99}  % results/attention/x2_breakdown.json
\newcommand{\FixedVonePct}{0.52}  % results/attention/x2_breakdown.json
\newcommand{\FixedVtwoUs}{15}  % results/attention/x2_breakdown.json
\newcommand{\FixedVtwoPct}{0.82}  % results/attention/x2_breakdown.json
\newcommand{\FixedVthreeUs}{2}  % results/attention/x2_breakdown.json
\newcommand{\FixedVthreePct}{0.11}  % results/attention/x2_breakdown.json
\newcommand{\NxNFloorMs}{1.07}  % results/attention/x2_breakdown.json
\newcommand{\NxNMB}{201}  % results/attention/x2_breakdown.json
\newcommand{\DepSavingMs}{0.40}  % results/attention/x2_breakdown.json
\newcommand{\CopyInMs}{1.42}  % results/attention/x2_breakdown.json
\newcommand{\CopyOutMs}{0.47}  % results/attention/x2_breakdown.json
\newcommand{\TierALower}{213}  % results/tokens/era_tra/summary.json
\newcommand{\TierASettings}{412}  % results/tokens/era_tra/summary.json
\newcommand{\TierBLower}{30}  % results/tokens/era_tra/summary.json
\newcommand{\TierBSettings}{412}  % results/tokens/era_tra/summary.json
\newcommand{\TierCLower}{3}  % results/tokens/era_tra/summary.json
\newcommand{\TierCSettings}{360}  % results/tokens/era_tra/summary.json
\newcommand{\CapClips}{100}  % results/gate_captions/scores.json
\newcommand{\CapDenseCider}{0.267}  % results/gate_captions/scores.json
\newcommand{\CapShuffledDelta}{0.218}  % results/gate_captions/scores.json
\newcommand{\CapBeightDelta}{\ensuremath{+0.001}}  % results/gate_captions/scores.json
\newcommand{\CapBeightLo}{\ensuremath{-0.038}}  % results/gate_captions/scores.json
\newcommand{\CapBeightHi}{\ensuremath{+0.040}}  % results/gate_captions/scores.json
\newcommand{\CapBeightErr}{27}  % results/gate_captions/scores.json
\newcommand{\CapCtwoDelta}{\ensuremath{+0.000}}  % results/gate_captions/scores.json
\newcommand{\CapCtwoLo}{\ensuremath{-0.041}}  % results/gate_captions/scores.json
\newcommand{\CapCtwoHi}{\ensuremath{+0.041}}  % results/gate_captions/scores.json
\newcommand{\CapCtwoErr}{21}  % results/gate_captions/scores.json
\newcommand{\CapHoldoneDelta}{\ensuremath{+0.003}}  % results/gate_captions/scores.json
\newcommand{\CapHoldoneLo}{\ensuremath{-0.039}}  % results/gate_captions/scores.json
\newcommand{\CapHoldoneHi}{\ensuremath{+0.043}}  % results/gate_captions/scores.json
\newcommand{\CapHoldoneErr}{36}  % results/gate_captions/scores.json
\newcommand{\CapBfloatSame}{65}  % results/gate_captions/scores.json


% Outline notes and figure placeholders, removed as the draft fills in.
\newcommand{\pending}[1]{{\color{red}[pending: #1]}}
\newcommand{\placeholder}[3]{\fbox{\parbox[c][#2][c]{\dimexpr#1-2\fboxsep-2\fboxrule\relax}%
  {\centering\footnotesize\color{red}#3}}}

\begin{document}
% IEEE reference style: more than six authors become "first author et al."
\bstctlcite{IEEEexample:BSTcontrol}

\title{Where the Time Goes in Attention on a 4\,GB Laptop GPU, and Why Temporal Gating Does Not Pay for an
Aerial Video Captioner\\{\color{red}\normalsize[working title]}}

\author{\IEEEauthorblockN{Ahad Ulla Baig}
\IEEEauthorblockA{\textit{Dept. of Artificial Intelligence} \\
\textit{and Machine Learning} \\
\textit{Ramaiah Institute of Technology}\\
Bengaluru, India \\
ahadullabaig.16@gmail.com}
\and
\IEEEauthorblockN{Divisha Gera}
\IEEEauthorblockA{\textit{Dept. of Artificial Intelligence} \\
\textit{and Machine Learning} \\
\textit{Ramaiah Institute of Technology}\\
Bengaluru, India \\
divisha.018@gmail.com}
\and
\IEEEauthorblockN{Charithra G}
\IEEEauthorblockA{\textit{Dept. of Artificial Intelligence} \\
\textit{and Machine Learning} \\
\textit{Ramaiah Institute of Technology}\\
Bengaluru, India \\
charithra754@gmail.com}
}

\maketitle

\begin{abstract}
\pending{written last. One sentence each: the setting (an offline video captioner whose vision encoder runs
attention at N = 1024, d = 64, 12 heads on an RTX 3050 Laptop GPU); the kernel ladder and its headline times at a
pinned \PinMhz\,MHz (\PinVoneMs, \PinVtwoMs, \PinVthreeMs\,ms); where the time goes (fixed costs at most
\FixedVtwoPct\%, fusion recovers \PinSavingPctOfNxN\% of the N$\times$N traffic it removes at equal clock); the
gating result (no tested gate beats dense at a lower frame rate at equal cost; captions show no detected
change).}
\end{abstract}

\begin{IEEEkeywords}
\pending{e.g. attention, CUDA, GPU kernels, kernel fusion, vision transformers, temporal redundancy, video
captioning}
\end{IEEEkeywords}

\section{Introduction}
\pending{
(1) Setting: an offline captioner for aerial video (SmolVLM2-500M~\cite{marafioti2025smolvlm}) whose vision
encoder runs dense attention at N = 1024 on a 4\,GB laptop GPU.
(2) Questions: where does attention's time go at this real shape (launch, compute, memory traffic, clock)?
Does frame-to-frame redundancy in aerial video justify a gated attention kernel?
(3) Approach: a from-scratch ladder (naive, tiled, fused) with times predicted before any kernel existed;
a shared benchmark protocol; Nsight Compute; an equal-clock comparison; token- and caption-level gating
measurements against an equal-cost baseline (dense at a lower frame rate).
(4) Contributions as measured results, no novelty claims until the literature search is done.
}

\section{Background and Related Work}
\pending{after the literature search (PLAN.md section 3.2 leads). Attention kernels and fusion
\cite{dao2022flashattention,dao2023flashattention2,bikshandi2023case}; temporal redundancy in video models
\cite{dutson2023eventful,parger2022deltacnn,parger2022motiondeltacnn,habibian2021skip,choudhury2024rlt};
dispatch and latency non-linearity at ViT scale \cite{abdellatif2026dispatch,eliopoulos2024pruning}; on-device
VLMs \cite{yang2026ondevice}. Check whether Eventful Transformers already compares against a reduced-frame-rate
baseline before calling our equal-cost comparison new.}

\section{Setup and Method}
\pending{
A. Hardware and measured limits: RTX 3050 Laptop (16 SMs, 4\,GB); FP32 FMA \PeakTflops\,TFLOPS sustained
(about \PeakGhz\,GHz, at the software power cap), DRAM \DramGBs\,GB/s, launch \LaunchUs\,\textmu s back to back.
B. Workload: the encoder's attention shape. C. Benchmark protocol (shared harness: interleaved groups,
per-sample clocks and throttle reasons, repeat runs). D. Predictions written before any attention kernel.
E. Correctness: every version against a float64 reference on all test cases. F. Clocks: the headline set is
pinned at \PinMhz\,MHz (one clock, one driver); unpinned times are reported as deployed, with their clocks.}

\section{The Kernel Ladder}
\pending{V1 naive (three kernels, N$\times$N in DRAM); V2 tiled (shared-memory tiles, register blocking,
softmax reading the matrix once); V3 fused (online softmax, no N$\times$N storage). Tile shapes chosen by a
sweep with the selection rule fixed beforehand.}

\section{Results}
\subsection{Scaling and the memory wall}
\pending{X1 figure, redrawn at paper width. V1 and V2 fail between N = 24k and 32k as the N$\times$N model
predicts; V3 runs at N = 32,768 (\VthreeBigNMs\,ms, one head).}

\begin{figure*}[t]
\centering
\placeholder{\textwidth}{2.4in}{Fig.: X1, time and memory against N (plots/x1\_plot.py, redrawn at 7.16\,in;
notes move to the caption)}
\caption{\pending{X1 caption.}}
\label{fig:x1}
\end{figure*}

\subsection{Predicted against measured}
\begin{table}[t]
\centering
\caption{Attention at 12 heads, N = 1024, D = 64, FP32. Ranges written before any attention kernel existed
(at the \PeakGhz\,GHz FMA peak), the same formulas rescaled after the fact to the pinned clock, the time at that
clock, and the time as deployed at each version's own clock.}
\label{tab:ladder}
\footnotesize
\setlength{\tabcolsep}{3pt}
% Generated by paper/numbers.py. Do not edit.
\begin{tabular}{@{}lccccc@{}}
\toprule
 & \multicolumn{2}{c}{Predicted (ms)} & Pinned & \multicolumn{2}{c}{As deployed} \\
\cmidrule(lr){2-3}\cmidrule(l){5-6}
Version & before & at \PinMhz\,MHz & (ms) & (ms) & (GHz) \\
\midrule
V1 (naive) & \PredVoneLo--\PredVoneHi & \PinVonePredLo--\PinVonePredHi & \PinVoneMs & \DepVoneMs & \DepVoneGhz \\
V2 (tiled) & \PredVtwoLo--\PredVtwoHi & \PinVtwoPredLo--\PinVtwoPredHi & \PinVtwoMs & \DepVtwoMs & \DepVtwoGhz \\
V3 (fused) & \PredVthreeLo--\PredVthreeHi & \PinVthreePredLo--\PinVthreePredHi & \PinVthreeMs & \DepVthreeMs & \DepVthreeGhz \\
B1 math & \PredBoneLo--\PredBoneHi & \PinBmathPredLo--\PinBmathPredHi & \PinBmathMs & \DepBmathMs & \DepBmathGhz \\
B1 mem.-eff. & \PredBoneLo--\PredBoneHi & \PinBeffPredLo--\PinBeffPredHi & \PinBeffMs & \DepBeffMs & \DepBeffGhz \\
\bottomrule
\end{tabular}

\end{table}
\pending{V2 lands inside its rescaled range, V3 is above, V1 is far above (it did not reuse operand rows from
L1 as the model assumed). The rescaled ranges are post hoc; say so.}

\subsection{Where the time goes}
\pending{X2: fixed costs \FixedVoneUs, \FixedVtwoUs, \FixedVthreeUs\,\textmu s (at most \FixedVtwoPct\% of a
version's time): H1's first clause holds. V2's N$\times$N traffic is worth \NxNFloorMs\,ms at full bandwidth;
V3 saves \PinSavingMs\,ms over V2 at equal clock (\DepSavingMs\,ms as deployed): the second clause fails.
Stall shares are descriptive; occupancy is untested for V3; the H1 criteria were written after timing.}

\begin{figure*}[t]
\centering
\placeholder{\textwidth}{2.2in}{Fig.: X2, stacked phases with floors (plots/x2\_plot.py), redrawn at paper
width; stall shares beside or below}
\caption{\pending{X2 caption.}}
\label{fig:x2}
\end{figure*}

\subsection{Against library attention}
\pending{FP32 only: at the pinned clock V3 \PinVthreeMs\,ms, PyTorch memory-efficient \PinBeffMs\,ms, math
\PinBmathMs\,ms. The demo runs bf16 and PyTorch's flash and cuDNN backends have no FP32 kernel: either add a
half-precision B1 or state the limitation here.}

\section{Does Temporal Gating Pay?}
\pending{
A. Pixel-level redundancy in ERA aerial video \cite{mou2020era} (X0).
B. Token traces: against dense at a lower frame rate with the output held, at equal cost, stale rows have the
lower connector error in \TierBLower{} of \TierBSettings{} settings, token deltas with exact attention in
\TierCLower{} of \TierCSettings, stale input (a floor no kernel reaches) in \TierALower{} of \TierASettings.
C. Captions (CapERA \cite{bashmal2023capera}, \CapClips{} test clips): no condition changes CIDEr-D
detectably; e.g.\ stale rows at threshold 8 (\CapBeightErr\% connector error)
\CapBeightDelta{} [\CapBeightLo, \CapBeightHi]. Resolution: another clip's references cost
\CapShuffledDelta. Held output at 1\,fps is near null by construction.
D. What it means for the kernel: the pipeline does not need high-rate encoding, so a gated kernel has no
application case here.}

\begin{figure*}[t]
\centering
\placeholder{\textwidth}{2.2in}{Fig.: token-trace error against cost per footage group
(plots/tokens\_era\_plot.py), redrawn at paper width}
\caption{\pending{Token-trace caption.}}
\label{fig:tokens}
\end{figure*}

\begin{figure}[t]
\centering
\placeholder{\columnwidth}{2.0in}{Fig.: caption quality against connector error and cost
(plots/gate\_captions\_plot.py), one panel at column width}
\caption{\pending{Caption-check caption.}}
\label{fig:captions}
\end{figure}

\section{Discussion and Limitations}
\pending{one GPU and one power-capped laptop; FP32 only (no tensor cores); one model and one attention shape;
the clock as a lever (fusion's unpinned advantage is mostly clock headroom, consistent with but not tested as a
power effect); occupancy untested for V3; H1 criteria not blind; caption metric resolution; copies of Q, K and
V would cost more than V3 (\CopyInMs\,+\,\CopyOutMs\,ms) if attention were offloaded alone.}

\section{Conclusion}
\pending{written last.}

\bibliographystyle{IEEEtran}
\bibliography{refs}

\end{document}
